% Clinical Background section. Paste into your own thesis, or \input it.

\section{Anatomical variants}
\label{sec:variants}

Beyond dominance, several variants change the connectivity of the coronary tree rather than
only its dimensions.
They are described here because a method that characterizes coronary topology must represent
them rather than discard them, and because two of them occur often enough in this cohort to
affect population statistics.

\subsection{Absent left main}

In a minority of people the left main artery does not exist as a discrete trunk, and the LAD
and LCX arise instead from separate ostia in the left aortic sinus
\cite{devos2024coronary,obrien2007anatomy}.
The left tree then has two roots rather than one.

The variant is present in 11 of the 800 cases, or 1.4\%, and it is identifiable from two
independent signals that agree exactly: those cases carry no LM label, and their left
centerlines resolve into two components each with its own ostium
(\texttt{scripts/survey\_topology.py}).
The practical consequence is that no stage of the analysis may assume a single ostium per
side.

\subsection{Ramus intermedius}

When the left main trifurcates rather than bifurcates, the additional vessel arising between
the LAD and the LCX is the ramus intermedius, labeled IM in this dataset
\cite{obrien2007anatomy}.
It supplies territory that would otherwise belong to a diagonal or an obtuse marginal branch.

A meta-analysis over 46 cadaveric, angiographic and CCTA studies covering 25,602 hearts puts
the aggregate prevalence at 25.5\%; restricted to CCTA cohorts it is 26.5\%, with a study
range of 19\% to 68.8\% \cite{dai2025ramus}.
An IM label is present in 196 of the 800 ImageCAS-X cases, or 24.5\%, consistent with that
pooled CCTA estimate.
The cohort figure is also stable across dominance groups, at 24.7\%, 22.0\% and 23.3\% for
right-, left- and codominant cases respectively.
The wide study range in the pooled estimate reflects both genuine population differences and
the threshold at which an observer records a third branch as a ramus intermedius rather than a
proximal diagonal \cite{dai2025ramus}.

\subsection{Myocardial bridging and anomalous origin}

A segment of a coronary artery may run within the myocardium rather than on its surface, a
variant termed myocardial bridging, most often affecting the mid-LAD
\cite{obrien2007anatomy}.
The bridged segment is compressed during systole.
Separately, a coronary artery may arise from the wrong sinus or from the pulmonary artery, a
group of anomalies whose clinical significance depends on whether the anomalous vessel takes
a course between the aorta and the pulmonary trunk \cite{obrien2007anatomy}.
Neither variant is separately labeled in this dataset, so neither can be studied directly
from the released annotations; both are noted because they are plausible explanations for a
coronary tree whose geometry departs markedly from the rest of a cohort.
